\section{Introduction}
A finite-dimensional orbit may admit a small positive factorization at each time and still require a growing register when those factorizations must be executed by one causal machine. For a fixed alphabet, that growth also depends on the distribution of its finite words. A qualitative assertion that all spectral gaps vanish loses precisely the arithmetic information which distinguishes such examples.

We study the numerical response
\begin{equation}\label{eq:target}
 \Pp(B=b\mid x,w,j)=\frac{1+b\rho\ip{e_j}{U_wx}}2,
 \qquad b\in\{-1,1\},
\end{equation}
where the command word $w$ has prescribed length $N$, the seed $x$ is a unit vector, and only one final coordinate $j$ is queried. The matrices in $U_w$ are orthogonal. A classical register must preserve these probabilities, or approximate every one of them within a fixed binary total-variation error. Its stochastic rows may vary arbitrarily with the epoch. Section~\ref{sec:model} defines the resource without a uniformity or finite-random-bit assumption.

The first main result is a pair of operational bounds in terms of finite data. For words of length $B$, let $\Gamma_{\calA}(k,B)$ be the largest, over word distributions, of the smallest spherical distortion after replacing the word orbit by $k$ centroid directions. This definition optimizes the converse law and permits every initial direction. Theorem~\ref{thm:profile} charges $-\log(1-\Gamma_{\calA}(K_t,B))$ to any interval of length $B$ ending at time $t$, even when all other registers are unrestricted. A finite interval recurrence combines all such charges and bounds actual final error. In the opposite direction, Theorem~\ref{thm:enclosure} turns a finite polytope and its worst command contraction into common stochastic rows. These are distinct lower and upper profiles; we do not assume that an optimal hidden realization must be a vector-state polytope.

For circle actions, Theorem~\ref{thm:circle-socp} removes the continuous search over centroid directions. The exact optimized packet value is a finite second-order cone problem indexed by cyclic partitions of the word phases. Thus a finite packet certificate is computationally different from an infinite harmonic-gap assertion. The theorem is not an efficiency assertion in a succinct word length.

Our arithmetic conclusions go beyond uniform bad approximation. For $\alpha\in\R^r$ with $1,\alpha_1,\ldots,\alpha_r$ rationally independent, write
\begin{equation}\label{eq:type-intro}
 \omega^*(\alpha)=\limsup_{\norm m_1\to\infty}
       \frac{-\log\norm{m\cdot\alpha}_{\R/\Z}}{\log\norm m_1}.
\end{equation}
The alphabet consists of the identity and the rotations through $2\pi\alpha_i$. Throughout, $r$, the alphabet, $0<\rho\le1/10$ and
$0\le\varepsilon<\rho/(2\sqrt2)$ are fixed.

\begin{theorem}[Minimal Diophantine type]\label{thm:metric-main}
If $\omega^*(\alpha)=r$, then
\begin{equation}\label{eq:log-law}
 \lim_{N\to\infty}\frac{\log W_{N,\varepsilon}(\alpha)}{\log N}
                         =\frac{r}{2r+1}.
\end{equation}
For Lebesgue-almost every $\alpha$, and every fixed $\delta>0$, there are positive constants $c,C$ such that, for $N\ge2$,
\begin{equation}\label{eq:ae-log-bounds}
 c\frac{N^{r/(2r+1)}}{(\log(N+2))^{2r(1+\delta)/(2r+1)}}
 \le W_{N,\varepsilon}(\alpha)
 \le C N^{r/(2r+1)}(\log(N+2))^{2r^2(1+\delta)/(2r+1)}.
\end{equation}
The upper inequalities are achieved by exact common-row machines.
\end{theorem}

The exponent in \eqref{eq:log-law} need not imply a two-sided constant-factor law. The earlier dual badly approximable theorem supplies the stronger $\Theta$ conclusion on its stated class, and is recovered in Corollary~\ref{cor:ba}. Theorem~\ref{thm:metric-main} instead allows arbitrarily small Diophantine constants along subsequences. The quantitative arithmetic ingredients are classical; the memory conclusions concern their compatible realization and all-hidden-state converse.

There is also a genuine failure of a single exponent, even for one fixed irrational angle.
\begin{theorem}[Irrational horizons and Liouville fluctuations]\label{thm:liouville-main}
For every irrational $\alpha$ there are horizons $M_n\to\infty$ such that
$W_{M_n,\varepsilon}(\alpha)=\Theta(M_n^{1/3})$. If $\alpha$ is Liouville, then
\begin{equation}\label{eq:liouville-summary}
 \liminf_{N\to\infty}\frac{\log W_{N,\varepsilon}(\alpha)}{\log N}=0,
 \qquad
 \frac13\le\limsup_{N\to\infty}\frac{\log W_{N,\varepsilon}(\alpha)}{\log N}
 \le\frac12.
\end{equation}
Nevertheless $W_{N,\varepsilon}(\alpha)\to\infty$. All displayed upper bounds have exact realizations.
\end{theorem}

The upper limit in \eqref{eq:liouville-summary} is not determined here. The result establishes oscillation between subpower subsequences and a positive-power subsequence; it does not replace that distinction by a guessed exponent for all irrationals.

\paragraph{Relation to previous work.}
Finite codebooks, nearest-center partitions and positive convex realizations are classical tools \cite{GL,Vidyasagar}. The endpoint conditional-centroid contraction used below was established in the preceding packet development \cite{GTF35,GTF42}; we give its complete proof and use it to define a finite alphabet profile. The general interval optimization, the circle word-law program, and the metric and Liouville memory conclusions are the new statements of this revision. Classical Dirichlet approximation, continued fractions and algebraic examples are not counted as new arithmetic. The direct comparison with Dick--Goda--Larcher--Pillichshammer--Suzuki \cite{DGLPS}, including their Theorem~3.9 and Lemma~3.11, is made in Section~\ref{sec:comparison}.

\begin{center}\small
\begin{tabular}{@{}p{.22\textwidth}p{.70\textwidth}@{}}\toprule
Quantity & Meaning \\\midrule
$K_t$, $W_{N,\varepsilon}$ & Available labels at a cut; minimum whole-run peak.\\
$\Gamma_{\calA}(k,B)$ & Best finite-word distortion against $k$ endpoint directions.\\
$\mathcal I_N(K_\bullet)$ & Maximum sum of distortion charges over disjoint intervals.\\
$\lambda_{\calA}(k,a)$ & Least common enclosure dilation with $k$ vertices and inradius $a$.\\
$\psi_\alpha(H)$, $\eta_\alpha(q)$ & Dual separation and simultaneous rational-approximation errors.\\\bottomrule
\end{tabular}
\end{center}
